\documentclass[aspectratio=169,11pt]{beamer}
\usetheme[progressbar=frametitle,numbering=fraction]{metropolis}
\usepackage[polish]{babel}
\usepackage{booktabs,tabularx}
\usepackage{amsmath}
\graphicspath{{../report/figures/}}

\definecolor{accent}{HTML}{4477AA}
\setbeamercolor{progress bar}{fg=accent}
\setbeamercolor{frametitle}{bg=black!85}
\setbeamercolor{alerted text}{fg=accent!80!black}

\title{Analiza miar różnorodności próbek w~baggingu}
\subtitle{Uczenie Zespołowe, projekt 1}
\author{Jerzy Czarnecki \and Małgorzata Dżaman}
\date{24.11.2026}
\institute{Politechnika Warszawska, Wydział MiNI}

\begin{document}

\maketitle

% --------------------------------------------------------------
\begin{frame}{Problem}
\small
\textbf{Bagging}: każde drzewo uczy się na osobnej próbce bootstrapowej (losowanie ze zwracaniem $n$ obserwacji z~$n$), predykcje drzew uśredniamy. \textbf{Las losowy}: bagging z~dodatkowym losowaniem cech w~każdym węźle.\\[0.4em]
\normalsize
\textbf{Jakość zespołu} zależy od siły drzew $s$ i~ich korelacji $\bar\rho$ \;(Breiman 2001):
\[
PE^{*} \le \frac{\bar{\rho}\,(1-s^{2})}{s^{2}}
\]
\vspace{-0.5em}
\begin{itemize}
  \item obie wielkości mierzymy zwykle \alert{po} uczeniu (predykcje, błędy),
  \item ale obie zależą od \alert{próbek}, które znamy \alert{przed} uczeniem.
\end{itemize}

\vspace{0.6em}
\begin{block}{Pytanie badawcze}
Czy miary różnorodności i~reprezentatywności próbek, obliczone przed uczeniem, są związane z~jakością zespołu, także po uwzględnieniu siły drzew?
\end{block}
\end{frame}

% --------------------------------------------------------------
\begin{frame}{Hipotezy}
\begin{description}[H1]
  \item[H1] Bardziej różne próbki $\Rightarrow$ bardziej różne predykcje drzew;\\ bardziej reprezentatywna próbka $\Rightarrow$ lepsze drzewo.
  \item[H2] Miary próbek wiążą się z~jakością zespołu także przy kontroli jakości drzew; przy stałym rozmiarze większe pokrywanie się próbek szkodzi.
  \item[H3] Co najmniej jeden sposób losowania daje inny wynik niż losowanie bootstrapowe.
\end{description}
\vspace{0.5em}
Wszystkie testy dwustronne, $\alpha = 0.05$, korekta Holma w~obrębie hipotezy.
\end{frame}

% --------------------------------------------------------------
\begin{frame}{15 wariantów losowania próbek ($M = 100$ drzew)}
\small
\begin{tabularx}{\textwidth}{lXX}
\toprule
& \textbf{jak losujemy} & \textbf{co zmienia} \\
\midrule
\textbf{S1} bootstrap & ze zwracaniem, $m = rn$, $r \in \{0.2,\dots,1\}$ & rozmiar próbki (siła i~różnorodność naraz) \\
\textbf{S2} bez zwracania & $m = sn$, $s = 1-e^{-r}$ & to samo $U$ co S1, bez powtórzeń \\
\textbf{S3} wspólny rdzeń & $m = 0.5n$; $\gamma m$ obserwacji \alert{wspólnych dla wszystkich} drzew, reszta osobno & \alert{tylko pokrywanie się} próbek \\
\textbf{S4} warstwowy & bootstrap osobno w~każdej klasie / przedziale $y$ & usuwa wahania proporcji klas \\
\bottomrule
\end{tabularx}

\vspace{0.6em}
Wariant odniesienia: S1~$r = 1$ (losowanie bootstrapowe, $m = n$). Wszystkie warianty uczymy na tych samych podziałach i~z~tymi samymi ziarnami.
\end{frame}

% --------------------------------------------------------------
\begin{frame}{Miary próbek i~wielkości po uczeniu}
\begin{columns}[T]
\begin{column}{0.5\textwidth}
\textbf{Przed uczeniem} (indeksy próbek i~część treningowa)\\[0.3em]
\emph{pary próbek} (4950 par):
\begin{itemize}
  \item $D_J$: odległość Jaccarda zbiorów obserwacji
  \item $D_C$: odległość z~uwzględnieniem powtórzeń
  \item $D_Y$: różnica rozkładów $y$
\end{itemize}
\emph{pojedyncza próbka} (100 drzew):
\begin{itemize}
  \item $U$: udział różnych obserwacji
  \item $R_Y$, $R_X$: odległość rozkładu $y$ i~cech od części treningowej
\end{itemize}
\end{column}
\begin{column}{0.5\textwidth}
\textbf{Po uczeniu} (część testowa)\\[0.3em]
\begin{itemize}
  \item jakość każdego drzewa \\ $\to$ \alert{siła}
  \item niezgodność par drzew lub $1-$korelacja błędów \\ $\to$ \alert{niezależność}
  \item zysk z~agregacji \\ (regresja: $=$ rozrzut predykcji, Krogh i~Vedelsby)
\end{itemize}
\end{column}
\end{columns}
\end{frame}

% --------------------------------------------------------------
\begin{frame}{Koncentracja miar próbek bootstrapowych}
\centering
\includegraphics[width=0.84\textwidth]{01_concentration.png}

\small
$P(\text{obserwacja w próbce}) \to 1-e^{-1} \approx 0.632$, \quad $D_J \to 1 - \tfrac{p}{2-p} \approx 0.538$, \quad rozrzut $\sim 1/\sqrt{n}$.
Wynika to z~własności losowania, dlatego różnorodność próbek zmieniamy \alert{celowo}, wybierając wariant losowania.
\end{frame}

% --------------------------------------------------------------
\begin{frame}{Dane i~protokół}
\begin{columns}[T]
\begin{column}{0.52\textwidth}
\begin{itemize}
  \item 12 zbiorów klasyfikacyjnych i~12 regresyjnych (208--20\,000 obserwacji, 4--81 cech)
  \item walidacja krzyżowa 5$\times$10, czyli 50 podziałów na zbiór
  \item las losowy z~domyślnymi parametrami \texttt{sklearn}, głosowanie miękkie; \\ w~klasyfikacji także bagging (bez losowania cech)
  \item 18\,000, 9\,000 i~2\,400 zespołów po 100 drzew (eksperyment główny, bagging, losowania)
\end{itemize}
\end{column}
\begin{column}{0.48\textwidth}
\textbf{Bez wycieku informacji}\\[0.2em]
\small
\emph{przed podziałem} tylko decyzje dotyczące pojedynczego wiersza lub kolumny:
usunięcie braków (bez imputacji), duplikatów, kolumn stałych\\[0.4em]
\emph{w~każdym podziale, na części treningowej}: one-hot, standaryzacja dla $R_X$, przedziały $y$ dla S4, losowanie próbek
\end{column}
\end{columns}
\end{frame}

% --------------------------------------------------------------
\begin{frame}{Statystyka: jednostką jest zbiór danych}
\begin{itemize}
  \item 50 podziałów nie jest niezależnych, więc wynikiem zbioru jest średnia lub mediana po podziałach
  \item H3: test Friedmana i~14 porównań z~wariantem odniesienia (Wilcoxon, Holm) (Demšar 2006)
  \item H1, H2: korelacje Spearmana w~zbiorze, potem test Wilcoxona po zbiorach; w~H2 także \alert{korelacja cząstkowa} z~kontrolą jakości drzew
\end{itemize}
\vspace{0.3em}
\begin{alertblock}{Mała liczba zbiorów, mała moc testów}
Przy 12 zbiorach najmniejsze możliwe $p$ to $0.0005$, a~próg Holma to $0.05/14 = 0.0036$, więc wariant musi być lepszy prawie we wszystkich zbiorach. Dlatego stosujemy też \textbf{analizę łączną 24 zbiorów}, \textbf{bayesowski test znakowanych rang} z~ROPE $\pm 1$\,p.p.\ (klasyfikacja) i~$\pm 1\%$ RMSE (regresja) oraz podajemy wielkości efektów.
\end{alertblock}
\end{frame}

% --------------------------------------------------------------
\begin{frame}{H3: żaden wariant nie jest lepszy od wariantu odniesienia}
\begin{columns}[T]
\begin{column}{0.6\textwidth}
\includegraphics[width=\textwidth]{05_average_ranks.png}
\end{column}
\begin{column}{0.4\textwidth}
\small
\begin{itemize}
  \item test Friedmana: $p<0.001$ w~obu zadaniach, więc sposób losowania \alert{wpływa} na wynik
  \item najlepsze S4 i~S2~$s{=}0.63$, ale tylko $+0.09$ i~$+0.04$\,p.p.; test bayesowski: $P(\text{ROPE}) = 1.00$
  \item \alert{istotnie gorsze}: wspólny rdzeń (S3, $\gamma \ge 0.25$), w~analizie 24 zbiorów także małe próbki
  \item stabilność: $\sigma/\sigma_0 \approx 1$, dla S3~$\gamma{=}0.75$ w~regresji $1.20$
\end{itemize}
\end{column}
\end{columns}
\end{frame}

% --------------------------------------------------------------
\begin{frame}{O~wyniku decyduje liczba różnych obserwacji i~pokrywanie się próbek}
\centering
\includegraphics[width=0.58\textwidth]{02_size_and_core.png}

\small
S1~$\approx$ S2 przy tym samym $U$, więc powtórzenia obserwacji nie wpływają na wynik. \quad
S3 (stałe $m$): $\gamma = 0$ lepsze od $\gamma = 0.75$ w~\alert{21 z~24} zbiorów, $p<0.001$.
\end{frame}

% --------------------------------------------------------------
\begin{frame}{Siła a~niezależność drzew}
\centering
\includegraphics[width=0.84\textwidth]{03_strength_vs_diversity.png}

\small
S1 i~S2 przesuwają się \alert{po krzywej} kompromisu: przy mniejszej próbce drzewa są bardziej różne, ale słabsze.\\
S3 zmienia \alert{tylko niezależność} przy stałej sile, więc bezpośrednio sprawdza oszacowanie Breimana.
\end{frame}

% --------------------------------------------------------------
\begin{frame}{H1 i~konkretne losowanie: bardzo słaby związek}
\begin{columns}[T]
\begin{column}{0.55\textwidth}
\includegraphics[width=\textwidth]{09_pair_example.png}
\end{column}
\begin{column}{0.45\textwidth}
\small
\begin{itemize}
  \item $D_J$ a~różnice predykcji par drzew: dodatnia korelacja w~10--12 z~12 zbiorów, istotna, ale \alert{$\rho \approx 0.02$--$0.03$}
  \item $U$, $R_Y$, $R_X$ a~jakość drzewa: brak związku
  \item 20 losowań tego samego zespołu: brak związku miar z~jakością ($|\tilde\rho| < 0.1$, wszystkie $p_H > 0.4$)
\end{itemize}
\vspace{0.3em}
Przyczyna: \alert{koncentracja} miar, prawie wszystkie pary mają $D_J \approx 0.54$.
\end{column}
\end{columns}
\end{frame}

% --------------------------------------------------------------
\begin{frame}{H2: przy tej samej sile drzew różnorodność próbek poprawia zespół}
\small
Korelacja miary z~jakością zespołu w~15 wariantach, mediana po zbiorach (w~nawiasie liczba zbiorów z~$\rho>0$)
\vspace{0.4em}

\centering
\begin{tabular}{lcccc}
\toprule
& \multicolumn{2}{c}{\textbf{prosta}} & \multicolumn{2}{c}{\textbf{cząstkowa} (kontrola siły drzew)} \\
\cmidrule(lr){2-3}\cmidrule(lr){4-5}
\textbf{Miara $D_J$} & $\tilde\rho$ & $p_H$ & $\tilde\rho$ & $p_H$ \\
\midrule
regresja & $-0.42$ (4/12) & 1.000 & $+0.59$ (\alert{12/12}) & \textbf{0.003} \\
klasyfikacja, las losowy & $-0.56$ (2/12) & 0.196 & $+0.28$ (9/12) & 0.205 \\
klasyfikacja, bagging & $-0.16$ (5/12) & 1.000 & $+0.54$ (\alert{12/12}) & \textbf{0.003} \\
\bottomrule
\end{tabular}

\raggedright
\vspace{0.3em}
\begin{itemize}
  \item bez kontroli kierunek zależy od zbioru: w~zbiorach z~dużym szumem (Adult, Abalone, Diabetes, Wine Quality) mniejsze próbki poprawiają wynik
  \item w~lesie losowym losowanie cech już zmniejsza korelację drzew, więc próbki wpływają słabiej
\end{itemize}
\end{frame}

% --------------------------------------------------------------
\begin{frame}{Pojedyncza miara nie wskazuje najlepszego wariantu}
\centering
\includegraphics[width=0.7\textwidth]{10_measures_vs_ranks.png}

\small
$U$ nie uwzględnia pokrywania się próbek (S3 przy $U=0.5$: rangi 6--14).
Związek $D_J$ z~rangą jest \alert{niemonotoniczny}: najgorsze warianty leżą na obu końcach.
\end{frame}

% --------------------------------------------------------------
\begin{frame}{Las losowy a~bagging}
\centering
\includegraphics[width=0.74\textwidth]{07_rf_vs_bagging.png}

\small
\begin{itemize}
  \item bagging: wspólny rdzeń $\gamma{=}0.75$ obniża macro-F1 o~$1.48$\,p.p.\ (las losowy: o~$0.48$\,p.p.)
  \item las losowy: drzewa bardziej różne w~12 z~12 zbiorów, ale większa wrażliwość na małe próbki\\
  Losowanie cech \alert{zastępuje} część różnorodności próbek.
\end{itemize}
\end{frame}

% --------------------------------------------------------------
\begin{frame}{Wnioski}
\footnotesize
\begin{itemize}
  \item \textbf{H1} potwierdzona częściowo: związek istotny, ale bardzo słaby, bo miary próbek są skoncentrowane
  \item \textbf{H2} potwierdzona przy kontroli siły drzew; pokrywanie się próbek pogarsza zespół (21 z~24 zbiorów)
  \item \textbf{H3} warianty się różnią, ale \alert{żaden nie jest lepszy} od wariantu odniesienia
\end{itemize}
\begin{block}{Wnioski praktyczne}
\begin{itemize}
  \item losowanie bootstrapowe jest dobrym wyborem domyślnym; o~wyniku decyduje $U$, a~nie zwracanie
  \item próbki różnych drzew nie powinny mieć wspólnej części
  \item rozmiar próbki warto dobierać dla zbioru (błąd \emph{out-of-bag}), zysk tylko przy dużym szumie
  \item miary obliczane przed uczeniem opisują sposób losowania, ale nie wskazują lepszego losowania
\end{itemize}
\end{block}
\end{frame}

% --------------------------------------------------------------
\begin{frame}[standout]
Dziękujemy!\\[0.5em]
\normalsize Pytania?
\end{frame}


\end{document}
